\subsection{A general limitation of fixed-target distribution tests}

\paragraph{D23.D1 (simultaneous distribution tests in a normally generated target).}
Let $F=\langle x,y\rangle$ be the free group, let $G$ be a finite
group generated by $X=q(x),Y=q(y)$, and assume that the normal
closure of $Y$ in $G$ is all of $G$. For $N\ge2$ let
\[
 K_N=\ker(F\longrightarrow\mathbb Z/N),\qquad
 x\longmapsto1,\quad y\longmapsto0.
\]
On the free basis $x^N,y_j=x^jyx^{-j}$, $0\le j<N$, define
\[
 \delta_N(x^N)=x,\quad\delta_N(y_0)=y,\quad
 \delta_N(y_j)=1\ (j>0),\qquad
 c_N(w)=\operatorname{exp}_y(\delta_N(w)).
\]
Given a finite set $\mathcal S\subset\{2,3,\ldots\}$, elements
$g_0,g_N\in G'$, and arbitrary integers $b_N$, there exists an
ordinary commutator word $f\in F'$ such that
\[
 q(f)=g_0,\qquad c_N(f)=b_N,\qquad
 q\delta_N(f)=Y^{b_N}g_N\quad(N\in\mathcal S).
\tag{D23.D1.1}
\]
For an infinite set $\mathcal S$ and arbitrary $b_N\in\widehat{\mathbb Z}$,
the same assertion holds with a profinite commutator word.

In particular, in one fixed such target the entire family of
Enriquez distribution equations can hold for a word with any
prescribed $G'$ value. This does not assert the gentle identities
for that word, or the distribution equations in other targets.

\begin{proof}
Write $L=\operatorname{lcm}(\mathcal S)$ and adjoin $L$ to
$\mathcal S$ if necessary, prescribing arbitrary extra data there.
The case of an empty set is immediate. For every positive $m$,
$q(K_m)=G$: this image is normal in $G$ and contains $Y$.
For $N\mid L$, the restriction $\delta_N(K_L)$ is exactly
$K_{L/N}$. This follows on the free basis: $x^L$ maps to
$x^{L/N}$, and $y_j$ maps to $x^{j/N}yx^{-j/N}$ if $N\mid j$
and to $1$ otherwise. Consequently
\[
 q\delta_N(K_L)=G.
\tag{D23.D1.2}
\]

First construct $b\in F'\cap K_L$ with $c_N(b)=b_N$ for all
the prescribed indices. Use the basis $x^L,y_0,\ldots,y_{L-1}$
of $K_L$. The coefficient of $y_0$ sets $c_L$ and affects every
$c_N$ by the same integer. In descending order of $N<L$ in
$\mathcal S$, the coefficient of $y_N$ sets the remaining
value of $c_N$ without altering a value already set at an index
$M>N$: indeed $c_M(y_N)=0$ for $M>N$, while $c_N(y_N)=1$.
Finally the coefficient of $y_1$ corrects the total $y$ exponent
to zero without altering any $c_N$. All these words have zero
$x$ exponent. The resulting word therefore belongs to $F'$.

Now put $V_L=K_L$. For $N<L$ put
\[
 V_N=\langle\!\langle y_N\rangle\!\rangle_{K_L}.
\]
By (D23.D1.2), the image of $V_N$ under $q\delta_N$ is the
normal closure of $XYX^{-1}$ in $G$, hence all of $G$.
For $V_L$ this follows directly from (D23.D1.2). Thus
\[
 q\delta_N(V_N')=G'.
\tag{D23.D1.3}
\]
If $M>N$ is another prescribed index, then $\delta_M(y_N)=1$
and hence $\delta_M(V_N)=1$. Every $V_N'$ lies in $F'$ and
has zero value under every homomorphism $c_M$.

Starting with $b$, proceed again in descending order of $N$.
The two elements $q\delta_N(b)$ and $Y^{b_N}g_N$ have the
same abelianization in $G$, since $\delta_N(b)$ has zero
$x$ exponent and $y$ exponent $b_N$. Their quotient therefore
belongs to $G'$. Equation (D23.D1.3) lets us correct it by
right multiplication by an element of $V_N'$. This leaves
all larger prescribed indices and all $c_M$ unchanged.
After finitely many corrections the distribution coordinates
have exactly the required values.

To set the original coordinate, put
$V_0=\langle\!\langle y_1\rangle\!\rangle_{K_L}$.
Its image under $q$ is all of $G$, because $q(K_L)=G$ and
$q(y_1)=XYX^{-1}$ normally generates $G$. Thus $q(V_0')=G'$.
But $\delta_N(V_0)=1$ for every $N\ge2$ dividing $L$.
One final factor from $V_0'$ adjusts $q(b)$ to $g_0$ without
changing any previously fixed coordinate. All constructions
are ordinary group operations with integer powers; finite
images of normal closures and their derived groups give
ordinary word representatives. This proves the finite assertion.

The same proof works in the profinite free group for finitely
many prescribed profinite integers, using their powers in the
profinite cyclic subgroups. Alternatively one can impose these
integers modulo successively larger moduli. All conditions in
(D23.D1.1) are closed conditions in the compact group
$\widehat F'$. Every finite subcollection has just been shown
to be satisfiable. Compactness therefore proves the infinite
assertion. It produces a single profinite word in this fixed
target, and does not promote that word to a gentle automorphism.
\end{proof}

\paragraph{D23.D2 (application to the compressed two-sign test).}
Let $G=E_0\times S_4$ be D23.P10's marked group of order
$5184$. Every one of its twelve necessary full-GT values
\[
 g_{\rho,e,\nu}=
 \left(\left(\binom{e}{\rho+(\rho+1)e},M_\rho\right),
               t^{2\nu}p^\rho\right),
 \quad\rho\in\mathbb F_3,\quad e,\nu\in\mathbb F_2,
\]
can be the value of one profinite word satisfying all the
Enriquez distribution equations in $G$, simultaneously for
all positive $N$, with arbitrarily prescribed $c_N$ for $N\ge2$.
In particular the fixed-target distribution equations do not
discard any of the six mixed-sign states. No actual gentle or
full-GT lift of these words is asserted.

\begin{proof}
The group is solvable. Its abelianization is $C_3\times C_2$,
and the image of $Y$ generates it: the affine linear part
$-L^2$ generates $\operatorname{SL}_2(\mathbb F_3)^{\rm ab}=C_3$,
and the four-cycle $t$ is odd. Hence the quotient of $G$ by
the normal closure of $Y$ is both solvable and perfect, so
is trivial. Every displayed state belongs to
$G'=(\mathbb F_3^2\rtimes Q_8)\times A_4$; this is also
checked in D23.P10's six-row certificate. Apply D23.D1 with
$g_0=g_N=g_{\rho,e,\nu}$ for every $N\ge2$. The case $N=1$
is the identity equation with $c_1=0$.
\end{proof}

\paragraph{Exact source and computational checks.}
The source theorem is B. Enriquez, \emph{Quasi-reflection algebras
and cyclotomic associators}, Selecta Math. 13 (2007),
Section 8, Theorem 8.1 and equation (50); arXiv:math/0408035v2,
printed pp.~30--32. It states these distribution equations
on the full profinite GT group. Remark 8.6 gives the precise
identity $c_{MN}=c_M+c_N$ under multiplication of the integer indices; that requirement
can be imposed on the arbitrary prescribed values above.
The supplied LM7\_PUSH E.6 already used this theorem in a
different finite target. The general normal-generation argument
here is a new proof of the stated fixed-target limitation.
As an independent finite check, \path{code/check_distribution_216.py}
constructs the induced wreath actions and verifies every one of
the twelve states for $N=2,3,4,6$. At each index the joint
projected commutator image has order $8957952=864^2\cdot12$;
all twelve choices of $c_N\bmod12$ occur. The proof above,
not those four finite checks, establishes the simultaneous
all-index statement.
